% 本文件由 LaTeX 排版 Agent 根据有效 translation.json 自动生成。
% author=Pamphilus; work_id=0615; title=《宗徒大事錄章旨闡釋》

\chapter{《宗徒大事錄章旨闡釋》}

% structure: p0001 source=h1 target=omitted reason=page-title-already-rendered-by-parent

我們既從父老和師長處獲益於方法與典範，遂謙卑地嘗試在本章旨闡釋中加以應用，懇請諸位寬恕我等之魯莽——無論就年齡或學識而言，我等尚屬年輕，並期望每位閱讀此篇文字者為我們祈禱，予以包容。因此，我們按照路加——福音作者兼歷史家——的記述，作出此闡釋。據此，我們以字母標示整章，並以星號標示各章的分段。\par

A. 論基督復活後的教導，及祂向門徒顯現，並應許賜下聖神之恩，以及基督升天的景象與方式。\par

B. 伯多祿向成為門徒者論述猶達斯的死亡與棄絕；* 本章亦包含藉祈禱及天主的恩寵以抽籤方式選出瑪弟亞以替代的段落。\par

C. 論五旬節聖神降臨於信徒身上。本章亦包含伯多祿的訓誨，* 及有關此事的先知言論，* 以及基督的苦難、復活、升天與聖神之恩；* 此外，在場者的信德，及他們藉聖洗所獲的救恩；* 再者，信徒間彼此團結的精神，促進公共利益，以及人數的增加。\par

D. 論藉基督之名醫治生來瘸腿者；及伯多祿的言論，他在其中辯論、同情並勸導關於救恩之事。此處有* 司祭長因嫉妒所發生的干預，及他們對奇蹟的判斷，以及伯多祿承認基督的能力與恩寵。另有段落論* 不信的司祭長禁止他們奉基督之名放膽講論，及宗徒被釋放。然後* 教會為宗徒忠信堅忍而獻上的感恩。\par

E. 論信徒間和諧而普世性的共融；* 亦論阿納尼雅與撒菲辣及其悲慘結局。\par

F. 論宗徒被投入監獄，夜間由主的天使領出，天使囑咐他們毫無阻礙地宣講耶穌；* 次日，司祭長再次逮捕他們，鞭打後斥責他們不可再教訓人。然後* 加瑪里耳關於宗徒的可靠意見，連同某些例子與證據。G. 論揀選七位執事。\par

H. 猶太人對斯德望的誣告，及斯德望關於天主與亞巴郎之盟約及十二位聖祖的演講。亦記述饑荒與買糧，雅各伯眾子彼此相認，以及梅瑟的出生和天主在西乃山向梅瑟顯現。* 亦論及以色列出埃及與鑄造金牛犢（及其他事件），直至撒羅滿時代及聖殿建造。* 然後論及耶穌基督超乎諸天的榮耀，這榮耀向斯德望本人啟示；為此緣故，斯德望被人用石頭砸死，並虔誠地安眠。\par

I. 論教會遭受迫害及斯德望之安葬；* 亦論宗徒斐理伯在撒瑪黎雅治好多人。\par

J. 論西滿·魔術士與許多人一同相信並受洗；* 亦論伯多祿與若望被派往他們那裡，為受洗者祈禱求聖神降臨。\par

K. 論聖神的恩賜不因金錢賜予，亦不賜予偽善者，而只賜予因信德而成聖的人；* 亦論西滿的偽善與斥責。\par

L. 論主在救恩道路上幫助善良而相信的人，正如從太監的事例所顯明。\par

M. 論從天上而來的神聖召叫，召喚保祿為基督的宗徒；* 亦論按天主的啟示，藉阿納尼雅之手醫治保祿並為他施洗，以及他放膽講論，並透過巴爾納伯與宗徒往來。\par

N. 論伯多祿在里達治好癱瘓的艾乃阿。* 亦記述寡婦之友塔彼達，伯多祿在約培藉祈禱使她從死者中復活。\par

O. 論科爾乃略，及天使對他所說的話。亦論從天上對伯多祿所說關於召叫外邦人的事。然後* 伯多祿被請而來到科爾乃略處。* 科爾乃略重述天使對他所說的事。* 伯多祿向他們教導基督，聽者領受聖神之恩，以及外邦信徒在當地受洗。\par

P. 論伯多祿向與他爭論的宗徒逐一敘述所有發生的事。然後論巴爾納伯被派往安提約基雅的弟兄們那裡。\par

Q. 阿加波關於天下饑荒的預言，以及送往耶路撒冷弟兄們的慷慨救濟。\par

R. 宗徒雅各伯被殺。* 亦論伯多祿被黑落德逮捕，及天使奉天命釋放他的方式，以及伯多祿夜間向門徒顯現後悄然退去。亦論守衛受罰，然後不虔敬的黑落德之悲慘而致命的覆滅。\par

S. 論巴爾納伯與保祿受聖神派遣前往塞浦路斯。* 論他們在那裡奉基督之名對厄里瑪斯術士所行的事。\par

T. 保祿關於基督真理的卓越闡述，依序引用法律與先知，兼涉歷史與福音；* 他使用駁斥與論證的方式，論及宣講之言轉向外邦人，以及他們所受的迫害與抵達依科尼雍。\par

U. 論他們在依科尼雍宣講基督時，許多人相信，宗徒卻遭受迫害。\par

V. 论在\\Place{吕斯特辣}生来瘸腿的人被宗徒治好；因此当地人以为他们是下凡的神明。此后，\\Person{保禄}在那里被邻近的人用石头砸。\par

W. 按宗徒的谕令和定断，信主的外邦人不应受割损。此处也有宗徒本人写给外邦人的信，论及他们应当远离的事。* \\Person{保禄}因\\Person{马尔谷}与\\Person{巴尔纳伯}起了争执。\par

X. 论\\Person{弟茂德}的教导，以及\\Person{保禄}按启示进入\\Place{马其顿}。* 论某位妇人\\Person{里狄雅}的信德和救恩，以及* 治好有占卜之神的使女，为此使女的主人将\\Person{保禄}下狱；以及* 在那里发生的地震和奇迹；以及狱警当即在同一夜全家信而受洗。* 宗徒们被恳求后从监牢出来。\par

Y. 论在\\Place{得撒洛尼}因他们的宣讲而起的骚乱，以及\\Person{保禄}逃往\\Place{贝洛雅}，又从那里到\\Place{雅典}。\par

Z. 论\\Place{雅典}祭坛上的题字，以及\\Person{保禄}的哲学宣讲和虔敬。\par

AA. 论\\Person{阿桂拉}和\\Person{普黎史拉}，以及\\Place{格林多}人的不信，以及天主按预知向\\Person{保禄}启示的善旨。也* 论会堂长\\Person{普黎斯苛}与某些人一起信而受洗。以及* 在\\Place{格林多}起了骚动，\\Person{保禄}离开；来到\\Place{厄弗所}，在那里讲论后便离去。* 论\\Person{阿颇罗}，一位有口才且信主的人。\par

BB. 论在\\Place{厄弗所}藉\\Person{保禄}祈祷，使信者领受圣洗和圣神的恩赐，以及治好民众。* 论\\Person{斯卡瓦}的众子，以及不可接近那些成为不信、不配于信德的人；以及信者的宣认；* 以及银匠\\Person{德默特琉}在\\Place{厄弗所}煽动反对宗徒的骚乱。\par

CC. 论\\Person{保禄}的巡游，其中也记载了在\\Place{特洛阿}的\\Person{欧提洛}死而藉祈祷复生；以及\\Person{保禄}在\\Place{厄弗所}对长老们的牧灵劝言；以及\\Person{保禄}从\\Place{厄弗所}到\\Place{凯撒勒雅}（\\Place{巴勒斯坦}）的航行。\par

DD. \\Person{阿加波}预言\\Person{保禄}在\\Place{耶路撒冷}将遭遇的事。\par

EE. \\Person{雅各伯}对\\Person{保禄}的嘱咐，关于他不应阻止希伯来人遵行割损。\par

FF. 论在\\Place{耶路撒冷}激起的反对\\Person{保禄}的骚乱，以及千夫长如何把他从暴民中救出。* 也\\Person{保禄}关于自己和蒙召为宗徒的辩护；* 以及\\Person{阿纳尼雅}在\\Place{大马士革}对\\Person{保禄}所说的话，以及一次在圣殿中临到他的异象和天主的声音。* 后来当\\Person{保禄}因这些话快要被鞭打时，他声明自己是罗马人，于是被释放。\par

GG. \\Person{保禄}下到公议会时，究竟受了什么苦，说了什么话，行了什么事。\par

HH. 论犹太人设伏谋害\\Person{保禄}，被\\Person{里息雅}发现；* 以及\\Person{保禄}被兵士和书信护送到\\Place{凯撒勒雅}的总督那里。\par

II. 论\\Person{特尔突罗}在\\Person{保禄}案中提出的控告，以及\\Person{保禄}在总督面前的自辩。\par

JJ. 论\\Person{斐理斯}被免职，\\Person{斐斯托}继任，以及\\Person{保禄}在他们面前的申辩和获释。\par

KK. \\Person{阿格黎帕}和\\Person{贝勒尼刻}到来，查问\\Person{保禄}的案件。* \\Person{保禄}在\\Person{阿格黎帕}和\\Person{贝勒尼刻}面前的自辩，关乎他自幼在律法中的教养，以及蒙召传福音。\\Person{阿格黎帕}对\\Person{斐斯托}说，\\Person{保禄}并未对犹太人做错什么。\par

LL. \\Person{保禄}往\\Place{罗马}的航行，充满许多极大危险。* \\Person{保禄}对同行者的劝言，论及他对得救的希望。\\Person{保禄}的船只失事，以及他们如何安全抵达\\Place{默里达}岛，并在岛上行了何等奇事。\par

MM. \\Person{保禄}如何从\\Place{默里达}抵达\\Place{罗马}。\par

NN. 论\\Person{保禄}在\\Place{罗马}与犹太人的谈话。\par

共计四十章；以下标有星号的章节共四十八段。\par

\SourceRecord{Translated by S.D.F. Salmond.；Ante-Nicene Fathers；Vol. 6.；Edited by Alexander Roberts, James Donaldson, and A. Cleveland Coxe.；Buffalo, NY: Christian Literature Publishing Co.,；1886.；<http://www.newadvent.org/fathers/0615.htm>.}
